\section{Related Work}
\label{sec:related_work}

\textbf{Regresson-based Object Detection Methods.}
Object detection has long been a cornerstone task in computer vision, with regression-based methods historically dominating the field. The core principle of these methods is to predict a bounding box by regressing its properties typically including center coordinates (x, y) and dimensions (width, height) as normalized offsets from a predefined reference.
Over the years, these methods have undergone significant evolution, progressing from early anchor-based CNN models like YOLO~\cite{redmon2016you}, SSD~\cite{liu2016ssd}, and Faster R-CNN~\cite{ren2016faster}, to anchor-free approaches such as CornerNet~\cite{law2018cornernet}, CenterNet~\cite{duan2019centernet}, and FCOS~\cite{tian2019fcos}. A major paradigm shift occurred with the introduction of Transformer-based detectors like DETR~\cite{carion2020end}, which framed object detection as a direct set prediction problem. This line of work was further advanced by models such as Deformable DETR~\cite{zhu2020deformable} and DINO~\cite{zhang2022dino}, which significantly improved performance and convergence speed.
Beyond these paradigmatic changes, the continuous improvement of regression-based detectors has been fueled by numerous incremental yet crucial innovations. These include architectural enhancements like Feature Pyramid Networks (FPN)\cite{lin2017feature}, advancements in loss functions like Focal Loss\cite{lin2017focal}, and sophisticated data augmentation techniques such as MixUp~\cite{zhang2017mixup} and Mosaic. It is the cumulative effect of these extensive and persistent efforts that has propelled regression-based object detectors to their current state of high performance and practical usability.

\textbf{Open-set Object Detection Methods.}
A long-term objective of object detection is to develop models capable of identifying an arbitrary number of object categories without task-specific fine-tuning, thereby addressing the challenges of real-world, dynamic scenarios. Open-set object detection represents a significant paradigm shift towards this goal, transcending the limitations of closed-set detection by empowering models to identify objects beyond a predefined set of categories. The prevalent approach to this challenge is text-prompted open-vocabulary object detection~\cite{li2022grounded, liu2023grounding, VLP:MDETR, yao2022detclip, zhong2022regionclip, ren2024dino, ghiasi2022scaling, cheng2024yolo, ma2023codet}. These methods typically leverage powerful pre-trained vision-language models like CLIP~\cite{VLP:CLIP} or BERT~\cite{devlin2018bert} to align textual descriptions with visual representations, demonstrating impressive zero-shot recognition capabilities. However, these models struggle with complex or nuanced descriptions due to their limited language understanding.
To overcome this, visual prompts~\cite{jiang2025t, jiang2023t, ren2024dino, wang2025yoloe, ravi2024sam, TransF:SAM, zou2023segment, li2024segment} have been introduced, allowing models to recognize objects using visual examples like boxes or points. Visual prompts are effective for rare or hard-to-describe objects but are less general than text prompts. Recent models like T-Rex2~\cite{jiang2025t} combine both text and visual prompts, using contrastive learning to leverage the strengths of each. This integration allows models to perform well across a wider range of object categories and real-world scenarios. While traditional open-set detectors achieve category-level generalization, they still lack deeper language understanding, making it challenging to handle context-rich real-world scenarios.

\textbf{MLLM-based Object Detection Methods.}
To overcome the shallow language understanding of traditional open-set detectors, a promising direction is to directly leverage the powerful reasoning capabilities of Multimodal Large Language Models (MLLMs) for object-level perception. The core idea is to reframe object detection as a language modeling task. Inspired by Pix2Seq~\cite{chen2021pix2seq}, a significant body of work has emerged that represents bounding box coordinates as a sequence of discrete, quantized tokens~\cite{peng2023kosmos, chen2023shikra, you2023ferret, wang2023cogvlm, zhan2025griffon}. These models, including Kosmos-2, Shikra, Ferret, and CogVLM, directly generate coordinate sequences through the standard next-token prediction mechanism of LLMs. This approach elegantly unifies object detection with the native capabilities of language models. However, as discussed in our introduction, this conceptually elegant approach faces significant practical challenges. While MLLMs excel at high-level image understanding, they often struggle with the fine-grained spatial precision required for object detection. Existing methods frequently suffer from limitations such as low recall rates, coordinate drift, and spurious duplicate predictions. We posit that these issues stem from two fundamental challenges: the inherent difficulty of learning a precise mapping from discrete tokens to a continuous pixel space using cross-entropy loss, and the behavioral deficiencies induced by the teacher-guided nature of Supervised Fine-Tuning (SFT). Addressing these challenges is the primary motivation for the design of Rex-Omni.
